\documentclass[12pt,a4paper]{article}

\usepackage[utf8]{inputenc}
\usepackage[T1]{fontenc}
\usepackage[indonesian]{babel}
\usepackage{amsmath,amssymb}
\usepackage{graphicx}
\usepackage{geometry}
\usepackage{setspace}
\usepackage{cite}
\usepackage{hyperref}

\geometry{margin=2.5cm}
\onehalfspacing

\title{\textbf{Spektroskopi Raman: Teori Dasar, Kaidah Seleksi, Raman Resonansi, dan Instrumentasi}}
\author{
  Anggota 1 \\
  Anggota 2 \\
  Anggota 3 \\
  \small Program Studi ... -- Mata Kuliah Spektroskopi (Semester 5)
}
\date{Desember 2026}

\begin{document}

\maketitle

\begin{abstract}
Tuliskan ringkasan makalah ilmiah di sini (maksimum 20 halaman total untuk seluruh dokumen, menggunakan ukuran font 12 pt sesuai ketentuan proyek kelompok).
\end{abstract}

\section{Pendahuluan}
Latar belakang singkat mengenai spektroskopi Raman dan tujuan penulisan makalah.

\section{Teori Dasar Hamburan dan Spektroskopi Raman}

% TODO: Isi pembahasan teori dasar hamburan Rayleigh, Stokes, Anti-Stokes,
% model polarisabilitas klasik, dan pendekatan kuantum.

\section{Perhitungan Kaidah Seleksi}

% TODO: Isi pembahasan perhitungan kaidah seleksi vibrasi dan rotasi Raman,
% serta analisis simetri/teori grup untuk menentukan mode aktif Raman.

\section{Teori Raman Resonansi}

% TODO: Isi pembahasan teori Raman resonansi, peningkatan intensitas hamburan,
% mekanisme Franck-Condon dan Herzberg-Teller, serta contoh aplikasinya.

\section{Perangkat Spektrometer Raman}

% TODO: Isi pembahasan komponen perangkat spektrometer Raman (sumber laser,
% sistem optik, filter Rayleigh, monokromator/spektrograf, dan detektor CCD).


\section{Kesimpulan}
Ringkasan poin-poin utama dari keempat topik yang telah dibahas.

\bibliographystyle{ieeetr}
\bibliography{references}

\end{document}
